\clearpage\thispagestyle{otFront}\phantomsection\label{references}\pdfbookmark[0]{Sources et édition}{bm-references}\otPageTitre{Sources et édition}{Références officielles · France}\fontsize{11}{15}\selectfont
\par\medskip\textbf{1. Arrêté et calendrier}\par Arrêté du 10 avril 2025 publié au BO n° 16 du 17 avril 2025, article 3 : le programme de mathématiques du cycle 3 s’applique en CM2 à la rentrée de l’année scolaire 2026-2027.\par\href{https://www.education.gouv.fr/bo/2025/Hebdo16/MENE2504620A}{\textcolor{otGreenDark}{Consulter la source officielle}}\par
\par\medskip\textbf{2. Programme officiel de mathématiques du cycle 3}\par Rubriques « Cours moyen deuxième année » : nombres, calcul et résolution de problèmes ; grandeurs et mesures ; espace et géométrie ; données et probabilités ; proportionnalité ; initiation à la pensée informatique.\par\href{https://www.education.gouv.fr/sites/default/files/programme-de-math-matiques-pour-le-cycle-3-439827.pdf}{\textcolor{otGreenDark}{Consulter la source officielle}}\par
\par\medskip\textbf{3. Accompagnement du programme}\par Éduscol : ressources d’accompagnement du programme de mathématiques au cycle 3, dont la résolution de problèmes en quatre phases et les schémas en barres.\par\href{https://eduscol.education.gouv.fr/5712/ressources-d-accompagnement-du-programme-de-mathematiques-au-cycle-3}{\textcolor{otGreenDark}{Consulter la source officielle}}\par
\par\bigskip Sources consultées le 23 octobre 2026.
Les cours sont rédigés pour ce livre ; les exercices, défis et bilans
sont originaux.\par\medskip
Périmètre du CM2 : nombres entiers jusqu’à 999 999 999 ; diviseurs
d’un nombre jusqu’à 100, liste complète jusqu’à 30, diviseurs et
multiples communs ; fractions de dénominateur au plus 60 et fractions
décimales jusqu’aux millièmes ; produit d’un entier par une fraction ;
divisions décimales par un diviseur à un chiffre ; programmes de calcul
jusqu’à trois instructions ; unités d’aire cm², dm², m² et aires du carré
et du rectangle ; angle droit de 90 degrés, sans rapporteur ; durées en
heures, minutes et secondes ; trapèzes, pentagones et hexagones ;
symétrie sur quadrillage, y compris le long d’une diagonale ; patrons du
cube et du pavé ; diagrammes circulaires ; probabilités « a chances
sur b », arbres et tableaux pour deux étapes.\par\medskip
La proportionnalité se traite par des raisonnements en langage naturel
(« trois fois plus, trois fois moins », addition de situations), sans
tableau de proportionnalité, sans coefficient ni produit en croix,
conformément au programme du cours moyen. Les défis approfondissent les
outils du livre et restent facultatifs. L’ordre des chapitres est une
proposition de progression.

\clearpage
